% Chapter 1 infographic -- use at the end of the chapter: % Chapter 1 infographic -- use at the end of the chapter: % Chapter 1 infographic -- use at the end of the chapter: % Chapter 1 infographic -- use at the end of the chapter: \input{chapters/ch01-infographic}
\ifcsname FEinfoLoaded\endcsname\else\input{common/fe-infographic}\fi
\begin{FEinfographic}{ch:01}{a long enough track record can reveal real skill}

\infopanel{1}{Three ways to make money in the market}{%
\infocards
  {Bearing risk}{Earning the equity premium, for example with a levered index fund, is a reward for risk, not evidence of skill.}
  {Luck}{Noise produces winners, such as coin-flip market timers, and we hear about the winners most.}
  {Skill}{A return beyond what the risk explains, after costs, from better information, analysis or execution.}}

\infopanel{2}{Evidence scoreboard}{%
\begin{infoscore}
\scoreitem{n}{No}{Mutual funds as a group, after fees}
\scoreitem{y}{Yes}{A minority of skilled managers}
\scoreitem{y}{Yes}{Records with an identifiable edge}
\scoreitem{n}{No}{Short records}
\scoreitem{n}{No}{The best of many backtests}
\end{infoscore}}

\infopanel{3}{Signal is small, noise is large: years of data needed to reach $t = 1.96$}{%
\infonumbers
  {61}{years for a Sharpe ratio of 0.25}
  {15}{years for a Sharpe ratio of 0.5}
  {4}{years for a Sharpe ratio of 1.0}
  {$t \approx \mathrm{SR}\sqrt{T}$}{evidence grows only with the square root of time}}

\infopanel{4}{Lab: 1,000 strategies with zero skill on the U.S. market since 1926}{%
\begin{infolab}
\item \textbf{Signal and noise.} The market's Sharpe ratio over the full sample and every 10-year window.
\item \textbf{Brilliant by chance.} About 25 random strategies look significant in the backtest, and the best reaches $t \approx 3.2$, yet the winners' later $t$ is about zero.
\item \textbf{The cost of searching.} The bar rises to $t > 3.89$ once 1,000 tries are allowed for.
\end{infolab}}

\infoverdict
  {A track record can reveal real skill, but only when it is long enough, clears a bar that allows for how many strategies were tried, and keeps working after it was chosen. Most records do not meet these conditions, and genuine skill appears to be rare.}
  {Judge evidence by how it was produced: ask how many strategies were tried, which part of the record is live, and how the strategy has done since its design was fixed.}
\end{FEinfographic}

\ifcsname FEinfoLoaded\endcsname\else% ======================================================================
% fe-infographic.tex -- shared layout for the one-page chapter infographics
% Loaded automatically by each chapters/chNN-infographic.tex.
% Works with the Glory style and the agentic style; colours are provided
% here if the preamble does not define them.
% ======================================================================
\providecolor{FEMaroon}{RGB}{128,30,45}
\providecolor{FECream}{RGB}{253,248,238}
\providecolor{FEGold}{RGB}{196,150,60}
\providecolor{fenavy}{HTML}{1B2A6B}
\providecolor{fegray}{HTML}{4D4D4D}
\providecolor{igRisk}{HTML}{1F4E79}
\providecolor{igLuck}{HTML}{8A8A8A}
\providecolor{igSkill}{HTML}{2E6B3F}
\providecolor{igYes}{HTML}{2E6B3F}
\providecolor{igNo}{HTML}{A33A3A}
\providecolor{igPart}{HTML}{C47F00}
\providecolor{igOpen}{HTML}{5D6D7E}

\def\FEinfoLoaded{1}

% Title band and page frame -----------------------------------------
\newenvironment{FEinfographic}[2]{% #1 = chapter label, #2 = claim
  \clearpage\thispagestyle{plain}%
  \begingroup\sffamily\footnotesize\setlength{\parskip}{0pt}\setlength{\parindent}{0pt}%
  \linespread{1.0}\selectfont
  \begin{tcolorbox}[enhanced, colback=FEMaroon, colframe=FEMaroon, arc=3pt,
      boxrule=0pt, left=8pt, right=8pt, top=5pt, bottom=5pt, before skip=0pt, after skip=5pt]
    {\color{white}\bfseries\large Chapter~\ref{#1} at a glance\par}
    \vspace{2pt}{\color{white}\small Claim on trial: \emph{#2}\par}
  \end{tcolorbox}}
  {\par\endgroup\clearpage}

% A numbered panel ----------------------------------------------------
\newcommand{\infopanel}[3]{% #1 number, #2 heading, #3 body
  \begin{tcolorbox}[enhanced, colback=FECream, colframe=FEGold, boxrule=0.6pt, arc=3pt,
      left=7pt, right=7pt, top=4pt, bottom=5pt, before skip=0pt, after skip=5pt]
    {\color{FEMaroon}\bfseries\small #1\quad #2\par}\vspace{3pt}%
    #3
  \end{tcolorbox}}

% Three idea cards ------------------------------------------------------
\newcommand{\infocard}[3]{% #1 colour, #2 title, #3 text
  \begin{minipage}[t]{0.315\linewidth}
    {\setlength{\fboxsep}{2.5pt}\colorbox{#1}{\makebox[\dimexpr\linewidth-2\fboxsep]{\color{white}\bfseries\rule[-0.3ex]{0pt}{2.2ex}#2}}}\par\vspace{3pt}
    {\scriptsize #3\par}
  \end{minipage}}
\newcommand{\infocards}[6]{%
  \infocard{igRisk}{#1}{#2}\hfill\infocard{igLuck}{#3}{#4}\hfill\infocard{igSkill}{#5}{#6}\par}

% Evidence scoreboard ---------------------------------------------------
\newcommand{\vchip}[2]{% #1 = y/n/p/o, #2 = label
  \edef\FEchipc{\ifx y#1igYes\else\ifx n#1igNo\else\ifx p#1igPart\else igOpen\fi\fi\fi}%
  \tcbox[on line, arc=2pt, boxrule=0pt, left=2pt, right=2pt, top=1pt, bottom=1pt,
    colback=\FEchipc,
    fontupper=\sffamily\bfseries\tiny\color{white}]{\makebox[3.6em]{#2}}}
% Items are laid out two to a row.
\newcount\FEscorecount
\newenvironment{infoscore}{\global\FEscorecount=0 \noindent\ignorespaces}{\par}
\newcommand{\scoreitem}[3]{% #1 class, #2 chip label, #3 situation
  \global\advance\FEscorecount by 1
  \begin{minipage}[t]{0.485\linewidth}\raggedright
    \vchip{#1}{#2}\hspace{4pt}\parbox[t]{\dimexpr\linewidth-5.4em}{\scriptsize #3}\vspace{3pt}
  \end{minipage}%
  \ifodd\FEscorecount\hfill\else\par\noindent\fi\ignorespaces}

% Four key numbers ------------------------------------------------------
\newcount\FEnumgroup
\newcommand{\infonum}[2]{%
  \begin{minipage}[t]{0.235\linewidth}
  \begin{tcolorbox}[enhanced, colback=white, colframe=FEGold, boxrule=0.5pt, arc=2pt,
      left=3pt, right=3pt, top=3pt, bottom=3pt, before skip=0pt, after skip=0pt,
      equal height group=FEnum\the\FEnumgroup, valign=top, halign=center]
    {\color{FEMaroon}\bfseries\Large #1\par}\vspace{2pt}{\scriptsize #2\par}
  \end{tcolorbox}
  \end{minipage}}
\newcommand{\infonumbers}[8]{%
  \global\advance\FEnumgroup by 1
  \infonum{#1}{#2}\hfill\infonum{#3}{#4}\hfill\infonum{#5}{#6}\hfill\infonum{#7}{#8}\par}

% Lab parts -------------------------------------------------------------
\newenvironment{infolab}{\scriptsize\begin{itemize}[leftmargin=1.1em, itemsep=1pt, topsep=1pt, parsep=0pt]}{\end{itemize}}

% Verdict band ----------------------------------------------------------
\newcommand{\infoverdict}[2]{% #1 = where the evidence stands, #2 = what investors can use
  \begin{tcolorbox}[enhanced, colback=white, colframe=FEMaroon, boxrule=0.8pt, arc=3pt,
      leftrule=4pt, left=7pt, right=7pt, top=4pt, bottom=5pt, before skip=0pt, after skip=0pt]
    {\color{FEMaroon}\bfseries\small Where the evidence stands\par}\vspace{2pt}
    {\scriptsize #1\par}\vspace{4pt}
    {\color{FEMaroon}\bfseries\small What investors can use\par}\vspace{2pt}
    {\scriptsize #2\par}
  \end{tcolorbox}}
\fi
\begin{FEinfographic}{ch:01}{a long enough track record can reveal real skill}

\infopanel{1}{Three ways to make money in the market}{%
\infocards
  {Bearing risk}{Earning the equity premium, for example with a levered index fund, is a reward for risk, not evidence of skill.}
  {Luck}{Noise produces winners, such as coin-flip market timers, and we hear about the winners most.}
  {Skill}{A return beyond what the risk explains, after costs, from better information, analysis or execution.}}

\infopanel{2}{Evidence scoreboard}{%
\begin{infoscore}
\scoreitem{n}{No}{Mutual funds as a group, after fees}
\scoreitem{y}{Yes}{A minority of skilled managers}
\scoreitem{y}{Yes}{Records with an identifiable edge}
\scoreitem{n}{No}{Short records}
\scoreitem{n}{No}{The best of many backtests}
\end{infoscore}}

\infopanel{3}{Signal is small, noise is large: years of data needed to reach $t = 1.96$}{%
\infonumbers
  {61}{years for a Sharpe ratio of 0.25}
  {15}{years for a Sharpe ratio of 0.5}
  {4}{years for a Sharpe ratio of 1.0}
  {$t \approx \mathrm{SR}\sqrt{T}$}{evidence grows only with the square root of time}}

\infopanel{4}{Lab: 1,000 strategies with zero skill on the U.S. market since 1926}{%
\begin{infolab}
\item \textbf{Signal and noise.} The market's Sharpe ratio over the full sample and every 10-year window.
\item \textbf{Brilliant by chance.} About 25 random strategies look significant in the backtest, and the best reaches $t \approx 3.2$, yet the winners' later $t$ is about zero.
\item \textbf{The cost of searching.} The bar rises to $t > 3.89$ once 1,000 tries are allowed for.
\end{infolab}}

\infoverdict
  {A track record can reveal real skill, but only when it is long enough, clears a bar that allows for how many strategies were tried, and keeps working after it was chosen. Most records do not meet these conditions, and genuine skill appears to be rare.}
  {Judge evidence by how it was produced: ask how many strategies were tried, which part of the record is live, and how the strategy has done since its design was fixed.}
\end{FEinfographic}

\ifcsname FEinfoLoaded\endcsname\else% ======================================================================
% fe-infographic.tex -- shared layout for the one-page chapter infographics
% Loaded automatically by each chapters/chNN-infographic.tex.
% Works with the Glory style and the agentic style; colours are provided
% here if the preamble does not define them.
% ======================================================================
\providecolor{FEMaroon}{RGB}{128,30,45}
\providecolor{FECream}{RGB}{253,248,238}
\providecolor{FEGold}{RGB}{196,150,60}
\providecolor{fenavy}{HTML}{1B2A6B}
\providecolor{fegray}{HTML}{4D4D4D}
\providecolor{igRisk}{HTML}{1F4E79}
\providecolor{igLuck}{HTML}{8A8A8A}
\providecolor{igSkill}{HTML}{2E6B3F}
\providecolor{igYes}{HTML}{2E6B3F}
\providecolor{igNo}{HTML}{A33A3A}
\providecolor{igPart}{HTML}{C47F00}
\providecolor{igOpen}{HTML}{5D6D7E}

\def\FEinfoLoaded{1}

% Title band and page frame -----------------------------------------
\newenvironment{FEinfographic}[2]{% #1 = chapter label, #2 = claim
  \clearpage\thispagestyle{plain}%
  \begingroup\sffamily\footnotesize\setlength{\parskip}{0pt}\setlength{\parindent}{0pt}%
  \linespread{1.0}\selectfont
  \begin{tcolorbox}[enhanced, colback=FEMaroon, colframe=FEMaroon, arc=3pt,
      boxrule=0pt, left=8pt, right=8pt, top=5pt, bottom=5pt, before skip=0pt, after skip=5pt]
    {\color{white}\bfseries\large Chapter~\ref{#1} at a glance\par}
    \vspace{2pt}{\color{white}\small Claim on trial: \emph{#2}\par}
  \end{tcolorbox}}
  {\par\endgroup\clearpage}

% A numbered panel ----------------------------------------------------
\newcommand{\infopanel}[3]{% #1 number, #2 heading, #3 body
  \begin{tcolorbox}[enhanced, colback=FECream, colframe=FEGold, boxrule=0.6pt, arc=3pt,
      left=7pt, right=7pt, top=4pt, bottom=5pt, before skip=0pt, after skip=5pt]
    {\color{FEMaroon}\bfseries\small #1\quad #2\par}\vspace{3pt}%
    #3
  \end{tcolorbox}}

% Three idea cards ------------------------------------------------------
\newcommand{\infocard}[3]{% #1 colour, #2 title, #3 text
  \begin{minipage}[t]{0.315\linewidth}
    {\setlength{\fboxsep}{2.5pt}\colorbox{#1}{\makebox[\dimexpr\linewidth-2\fboxsep]{\color{white}\bfseries\rule[-0.3ex]{0pt}{2.2ex}#2}}}\par\vspace{3pt}
    {\scriptsize #3\par}
  \end{minipage}}
\newcommand{\infocards}[6]{%
  \infocard{igRisk}{#1}{#2}\hfill\infocard{igLuck}{#3}{#4}\hfill\infocard{igSkill}{#5}{#6}\par}

% Evidence scoreboard ---------------------------------------------------
\newcommand{\vchip}[2]{% #1 = y/n/p/o, #2 = label
  \edef\FEchipc{\ifx y#1igYes\else\ifx n#1igNo\else\ifx p#1igPart\else igOpen\fi\fi\fi}%
  \tcbox[on line, arc=2pt, boxrule=0pt, left=2pt, right=2pt, top=1pt, bottom=1pt,
    colback=\FEchipc,
    fontupper=\sffamily\bfseries\tiny\color{white}]{\makebox[3.6em]{#2}}}
% Items are laid out two to a row.
\newcount\FEscorecount
\newenvironment{infoscore}{\global\FEscorecount=0 \noindent\ignorespaces}{\par}
\newcommand{\scoreitem}[3]{% #1 class, #2 chip label, #3 situation
  \global\advance\FEscorecount by 1
  \begin{minipage}[t]{0.485\linewidth}\raggedright
    \vchip{#1}{#2}\hspace{4pt}\parbox[t]{\dimexpr\linewidth-5.4em}{\scriptsize #3}\vspace{3pt}
  \end{minipage}%
  \ifodd\FEscorecount\hfill\else\par\noindent\fi\ignorespaces}

% Four key numbers ------------------------------------------------------
\newcount\FEnumgroup
\newcommand{\infonum}[2]{%
  \begin{minipage}[t]{0.235\linewidth}
  \begin{tcolorbox}[enhanced, colback=white, colframe=FEGold, boxrule=0.5pt, arc=2pt,
      left=3pt, right=3pt, top=3pt, bottom=3pt, before skip=0pt, after skip=0pt,
      equal height group=FEnum\the\FEnumgroup, valign=top, halign=center]
    {\color{FEMaroon}\bfseries\Large #1\par}\vspace{2pt}{\scriptsize #2\par}
  \end{tcolorbox}
  \end{minipage}}
\newcommand{\infonumbers}[8]{%
  \global\advance\FEnumgroup by 1
  \infonum{#1}{#2}\hfill\infonum{#3}{#4}\hfill\infonum{#5}{#6}\hfill\infonum{#7}{#8}\par}

% Lab parts -------------------------------------------------------------
\newenvironment{infolab}{\scriptsize\begin{itemize}[leftmargin=1.1em, itemsep=1pt, topsep=1pt, parsep=0pt]}{\end{itemize}}

% Verdict band ----------------------------------------------------------
\newcommand{\infoverdict}[2]{% #1 = where the evidence stands, #2 = what investors can use
  \begin{tcolorbox}[enhanced, colback=white, colframe=FEMaroon, boxrule=0.8pt, arc=3pt,
      leftrule=4pt, left=7pt, right=7pt, top=4pt, bottom=5pt, before skip=0pt, after skip=0pt]
    {\color{FEMaroon}\bfseries\small Where the evidence stands\par}\vspace{2pt}
    {\scriptsize #1\par}\vspace{4pt}
    {\color{FEMaroon}\bfseries\small What investors can use\par}\vspace{2pt}
    {\scriptsize #2\par}
  \end{tcolorbox}}
\fi
\begin{FEinfographic}{ch:01}{a long enough track record can reveal real skill}

\infopanel{1}{Three ways to make money in the market}{%
\infocards
  {Bearing risk}{Earning the equity premium, for example with a levered index fund, is a reward for risk, not evidence of skill.}
  {Luck}{Noise produces winners, such as coin-flip market timers, and we hear about the winners most.}
  {Skill}{A return beyond what the risk explains, after costs, from better information, analysis or execution.}}

\infopanel{2}{Evidence scoreboard}{%
\begin{infoscore}
\scoreitem{n}{No}{Mutual funds as a group, after fees}
\scoreitem{y}{Yes}{A minority of skilled managers}
\scoreitem{y}{Yes}{Records with an identifiable edge}
\scoreitem{n}{No}{Short records}
\scoreitem{n}{No}{The best of many backtests}
\end{infoscore}}

\infopanel{3}{Signal is small, noise is large: years of data needed to reach $t = 1.96$}{%
\infonumbers
  {61}{years for a Sharpe ratio of 0.25}
  {15}{years for a Sharpe ratio of 0.5}
  {4}{years for a Sharpe ratio of 1.0}
  {$t \approx \mathrm{SR}\sqrt{T}$}{evidence grows only with the square root of time}}

\infopanel{4}{Lab: 1,000 strategies with zero skill on the U.S. market since 1926}{%
\begin{infolab}
\item \textbf{Signal and noise.} The market's Sharpe ratio over the full sample and every 10-year window.
\item \textbf{Brilliant by chance.} About 25 random strategies look significant in the backtest, and the best reaches $t \approx 3.2$, yet the winners' later $t$ is about zero.
\item \textbf{The cost of searching.} The bar rises to $t > 3.89$ once 1,000 tries are allowed for.
\end{infolab}}

\infoverdict
  {A track record can reveal real skill, but only when it is long enough, clears a bar that allows for how many strategies were tried, and keeps working after it was chosen. Most records do not meet these conditions, and genuine skill appears to be rare.}
  {Judge evidence by how it was produced: ask how many strategies were tried, which part of the record is live, and how the strategy has done since its design was fixed.}
\end{FEinfographic}

\ifcsname FEinfoLoaded\endcsname\else% ======================================================================
% fe-infographic.tex -- shared layout for the one-page chapter infographics
% Loaded automatically by each chapters/chNN-infographic.tex.
% Works with the Glory style and the agentic style; colours are provided
% here if the preamble does not define them.
% ======================================================================
\providecolor{FEMaroon}{RGB}{128,30,45}
\providecolor{FECream}{RGB}{253,248,238}
\providecolor{FEGold}{RGB}{196,150,60}
\providecolor{fenavy}{HTML}{1B2A6B}
\providecolor{fegray}{HTML}{4D4D4D}
\providecolor{igRisk}{HTML}{1F4E79}
\providecolor{igLuck}{HTML}{8A8A8A}
\providecolor{igSkill}{HTML}{2E6B3F}
\providecolor{igYes}{HTML}{2E6B3F}
\providecolor{igNo}{HTML}{A33A3A}
\providecolor{igPart}{HTML}{C47F00}
\providecolor{igOpen}{HTML}{5D6D7E}

\def\FEinfoLoaded{1}

% Title band and page frame -----------------------------------------
\newenvironment{FEinfographic}[2]{% #1 = chapter label, #2 = claim
  \clearpage\thispagestyle{plain}%
  \begingroup\sffamily\footnotesize\setlength{\parskip}{0pt}\setlength{\parindent}{0pt}%
  \linespread{1.0}\selectfont
  \begin{tcolorbox}[enhanced, colback=FEMaroon, colframe=FEMaroon, arc=3pt,
      boxrule=0pt, left=8pt, right=8pt, top=5pt, bottom=5pt, before skip=0pt, after skip=5pt]
    {\color{white}\bfseries\large Chapter~\ref{#1} at a glance\par}
    \vspace{2pt}{\color{white}\small Claim on trial: \emph{#2}\par}
  \end{tcolorbox}}
  {\par\endgroup\clearpage}

% A numbered panel ----------------------------------------------------
\newcommand{\infopanel}[3]{% #1 number, #2 heading, #3 body
  \begin{tcolorbox}[enhanced, colback=FECream, colframe=FEGold, boxrule=0.6pt, arc=3pt,
      left=7pt, right=7pt, top=4pt, bottom=5pt, before skip=0pt, after skip=5pt]
    {\color{FEMaroon}\bfseries\small #1\quad #2\par}\vspace{3pt}%
    #3
  \end{tcolorbox}}

% Three idea cards ------------------------------------------------------
\newcommand{\infocard}[3]{% #1 colour, #2 title, #3 text
  \begin{minipage}[t]{0.315\linewidth}
    {\setlength{\fboxsep}{2.5pt}\colorbox{#1}{\makebox[\dimexpr\linewidth-2\fboxsep]{\color{white}\bfseries\rule[-0.3ex]{0pt}{2.2ex}#2}}}\par\vspace{3pt}
    {\scriptsize #3\par}
  \end{minipage}}
\newcommand{\infocards}[6]{%
  \infocard{igRisk}{#1}{#2}\hfill\infocard{igLuck}{#3}{#4}\hfill\infocard{igSkill}{#5}{#6}\par}

% Evidence scoreboard ---------------------------------------------------
\newcommand{\vchip}[2]{% #1 = y/n/p/o, #2 = label
  \edef\FEchipc{\ifx y#1igYes\else\ifx n#1igNo\else\ifx p#1igPart\else igOpen\fi\fi\fi}%
  \tcbox[on line, arc=2pt, boxrule=0pt, left=2pt, right=2pt, top=1pt, bottom=1pt,
    colback=\FEchipc,
    fontupper=\sffamily\bfseries\tiny\color{white}]{\makebox[3.6em]{#2}}}
% Items are laid out two to a row.
\newcount\FEscorecount
\newenvironment{infoscore}{\global\FEscorecount=0 \noindent\ignorespaces}{\par}
\newcommand{\scoreitem}[3]{% #1 class, #2 chip label, #3 situation
  \global\advance\FEscorecount by 1
  \begin{minipage}[t]{0.485\linewidth}\raggedright
    \vchip{#1}{#2}\hspace{4pt}\parbox[t]{\dimexpr\linewidth-5.4em}{\scriptsize #3}\vspace{3pt}
  \end{minipage}%
  \ifodd\FEscorecount\hfill\else\par\noindent\fi\ignorespaces}

% Four key numbers ------------------------------------------------------
\newcount\FEnumgroup
\newcommand{\infonum}[2]{%
  \begin{minipage}[t]{0.235\linewidth}
  \begin{tcolorbox}[enhanced, colback=white, colframe=FEGold, boxrule=0.5pt, arc=2pt,
      left=3pt, right=3pt, top=3pt, bottom=3pt, before skip=0pt, after skip=0pt,
      equal height group=FEnum\the\FEnumgroup, valign=top, halign=center]
    {\color{FEMaroon}\bfseries\Large #1\par}\vspace{2pt}{\scriptsize #2\par}
  \end{tcolorbox}
  \end{minipage}}
\newcommand{\infonumbers}[8]{%
  \global\advance\FEnumgroup by 1
  \infonum{#1}{#2}\hfill\infonum{#3}{#4}\hfill\infonum{#5}{#6}\hfill\infonum{#7}{#8}\par}

% Lab parts -------------------------------------------------------------
\newenvironment{infolab}{\scriptsize\begin{itemize}[leftmargin=1.1em, itemsep=1pt, topsep=1pt, parsep=0pt]}{\end{itemize}}

% Verdict band ----------------------------------------------------------
\newcommand{\infoverdict}[2]{% #1 = where the evidence stands, #2 = what investors can use
  \begin{tcolorbox}[enhanced, colback=white, colframe=FEMaroon, boxrule=0.8pt, arc=3pt,
      leftrule=4pt, left=7pt, right=7pt, top=4pt, bottom=5pt, before skip=0pt, after skip=0pt]
    {\color{FEMaroon}\bfseries\small Where the evidence stands\par}\vspace{2pt}
    {\scriptsize #1\par}\vspace{4pt}
    {\color{FEMaroon}\bfseries\small What investors can use\par}\vspace{2pt}
    {\scriptsize #2\par}
  \end{tcolorbox}}
\fi
\begin{FEinfographic}{ch:01}{a long enough track record can reveal real skill}

\infopanel{1}{Three ways to make money in the market}{%
\infocards
  {Bearing risk}{Earning the equity premium, for example with a levered index fund, is a reward for risk, not evidence of skill.}
  {Luck}{Noise produces winners, such as coin-flip market timers, and we hear about the winners most.}
  {Skill}{A return beyond what the risk explains, after costs, from better information, analysis or execution.}}

\infopanel{2}{Evidence scoreboard}{%
\begin{infoscore}
\scoreitem{n}{No}{Mutual funds as a group, after fees}
\scoreitem{y}{Yes}{A minority of skilled managers}
\scoreitem{y}{Yes}{Records with an identifiable edge}
\scoreitem{n}{No}{Short records}
\scoreitem{n}{No}{The best of many backtests}
\end{infoscore}}

\infopanel{3}{Signal is small, noise is large: years of data needed to reach $t = 1.96$}{%
\infonumbers
  {61}{years for a Sharpe ratio of 0.25}
  {15}{years for a Sharpe ratio of 0.5}
  {4}{years for a Sharpe ratio of 1.0}
  {$t \approx \mathrm{SR}\sqrt{T}$}{evidence grows only with the square root of time}}

\infopanel{4}{Lab: 1,000 strategies with zero skill on the U.S. market since 1926}{%
\begin{infolab}
\item \textbf{Signal and noise.} The market's Sharpe ratio over the full sample and every 10-year window.
\item \textbf{Brilliant by chance.} About 25 random strategies look significant in the backtest, and the best reaches $t \approx 3.2$, yet the winners' later $t$ is about zero.
\item \textbf{The cost of searching.} The bar rises to $t > 3.89$ once 1,000 tries are allowed for.
\end{infolab}}

\infoverdict
  {A track record can reveal real skill, but only when it is long enough, clears a bar that allows for how many strategies were tried, and keeps working after it was chosen. Most records do not meet these conditions, and genuine skill appears to be rare.}
  {Judge evidence by how it was produced: ask how many strategies were tried, which part of the record is live, and how the strategy has done since its design was fixed.}
\end{FEinfographic}
